\section{Related Work}
\label{sec:related-work}

\subsection{Software-Agent Benchmarks}

Code-agent evaluation has gradually moved from isolated function synthesis toward repository-level and tool-using software engineering. Function-level benchmarks such as HumanEval~\cite{chen2021codex} and MBPP~\cite{austin2021program} measure local programming ability under relatively self-contained problem statements~\cite{chen2021evaluating,austin2021programsynthesislargelanguage}. 
Later benchmarks extend the evaluation context from single functions to files, repositories, and executable development environments. 
RepoBench~\cite{liu2023repobench} evaluates repository-level code completion with cross-file retrieval and completion tasks, while M2RC-Eval~\cite{liu2025m2rc} studies multi-file and multi-repository contextual reasoning. 
Issue-resolution benchmarks further place agents inside existing repositories and ask them to generate patches for real software issues, including SWE-bench~\cite{jimenez2024swebench}, Multi-SWE-bench~\cite{zan2025multi}, SWE-Bench Pro~\cite{deng2025swe}, and SWE-Compass~\cite{xu2025swe}. 
Recent surveys summarize this broader transition from code models to autonomous software-engineering agents~\cite{llm_codegen_survey2024,llm_agents_se_survey2024,codegen_agents_survey2025}.
These benchmarks are essential for measuring local coding, repository context use, tool execution, and issue-driven editing. However, they mostly assume that the target software context already exists and that the agent`s goal is to complete, edit, or repair it. They therefore do not directly evaluate the emerging project-building setting, which is main focus of the following works as well as ours.
% While these benchmarks improve realism, they typically assume well-specified targets for existing projects. \rcb{} addresses a missing axis: evaluating an agent's ability to clarify ambiguous intents and construct complete repositories from scratch.

\subsection{0-to-1 Repository Generation}

Existing 0-to-1 repository benchmarks differ mainly in how much of the target project is visible to the agent before implementation begins. One common design is to preserve substantial implementation scaffolding while removing selected code only. Commit0~\cite{zhao2024commit0}, for example, keeps repository structure, class definitions, and function signatures, so the core challenge is to fill in missing implementations rather than clarify underspecified product requirements. A second design provides a full natural-language project specification and asks the agent to realize it in an empty workspace. NL2RepoBench~\cite{ding2026nl2repobench}, RepoGenesis~\cite{peng2026repogenesisbenchmarkingendtoendmicroservice}, PRDBench~\cite{fu2025prdbench}, ProjectEval~\cite{liu2025projecteval}, and RealBench~\cite{li2026realbench} all move closer to end-to-end project building, but they still largely assume that the agent has been given the decisive requirements needed to start implementation.

Another line of work relaxes implementation constraints while keeping the behavioral target available through static artifacts. ProgramBench~\cite{yang2026programbench} evaluates whether agents can reconstruct behaviorally equivalent programs from executables and documentation, emphasizing black-box behavioral fidelity. This is closer to our language-agnostic evaluation philosophy, but its target behavior is still directly specified rather than partially hidden behind an interactive clarification process.

Table~\ref{tab:benchmark-comparison} summarizes these differences along the dimensions most relevant to interactive project generation: project scale, language coverage, initially visible information, clarification interaction, and cross-language evaluation. Compared with existing 0-to-1 benchmarks, \rcb{} exposes less decisive information at the start, while allowing hidden requirements to be recovered through multi-turn interaction. It also evaluates generated repositories with language-agnostic black-box tests, enabling cross-language implementations. Thus, \rcb{} complements prior 0-to-1 benchmarks by moving the focus from implementing a given specification to discovering and implementing a partially specified product target.

%Taken together, these benchmarks increasingly evaluate long-horizon repository construction, dependency management, and behavioral correctness, but they mostly study what happens after the task specification is already available. Table~\ref{tab:benchmark-comparison} shows that \rcb{} differs from existing 0-to-1 benchmarks on every axis: with 480 projects across 12 mainstream languages, it is both the largest and the most multilingual of these benchmarks, and it is the only one to support clarification interaction and cross-language evaluation.

\subsection{Ambiguity and Clarification}

A growing line of work studies how code models handle unclear, incomplete, or interactive specifications~\cite{jin2026reqelicitgymevaluationenvironmentinterview,li2023pythoncodegenerationasking,Fakhoury_2024}. 
HumanEvalComm modifies HumanEval-style problems with inconsistency, ambiguity, and incompleteness to evaluate whether models ask useful clarification questions~\cite{Wu_2025}. 
Orchid studies the impact of requirement ambiguity on function-level code generation and categorizes ambiguity into lexical, syntactic, semantic, and vagueness-related forms~\cite{yang2026assessingimpactrequirementambiguity}. 
When Benchmarks Talk converts static code benchmarks into interactive settings where models must obtain key missing information from a simulated user~\cite{pan2025benchmarkstalkreevaluatingcode}. 
Ask or Assume evaluates clarification-seeking coding agents on an underspecified variant of SWE-bench Verified, studying whether agents can recognize when issue descriptions lack crucial context~\cite{edwards2026askassumeuncertaintyawareclarificationseeking}. 
These studies establish clarification as an important coding capability, but they mostly evaluate it in function-level tasks or issue-resolution settings. In such settings, clarification usually affects local behavior or a patch inside an existing repository. \rcb{} brings clarification into 0-to-1 repository generation, where missing requirements can affect architecture, public interfaces, dependencies, test protocols, and edge-case behavior. 
%It therefore jointly evaluates the clarification trajectory and the final executable repository, connecting the ambiguity-focused literature with the repository-generation benchmarks.
%While highlighting the importance of clarification, existing studies mostly focus on function-level or issue-resolution tasks. \rcb{} extends this to repository-level generation, jointly measuring implementation correctness and multi-turn interaction quality.